\section{Replication and Preservation: Complete Methods}
\label{app:newmethods}
\subsection{Study boundaries}
The starting study contains 72 matched-assignment fits.
Its full methods and results appear in Appendix~\ref{app:historicalv3}.
The nonlinear and RN50 replications each contain 45 fits.
Their protocols are \path{docs/STRENGTHEN_REPLICATION_PROTOCOL.json} and \path{docs/STRENGTHEN_ENCODER_REPLICATION_PROTOCOL.json}.
The frozen retention design contributes 180 fits under \path{results/strengthen_retention/protocol_v3.json}.
The preservation extension has its own protocol, \path{results/allocation_distillation/protocol_v1.json}.
This extension adds the allocation-plus-distillation comparison to the earlier design.
It uses previously examined test sets and is identified as exploratory.
The scientific record keeps every earlier protocol separately.
Some Source, U, A, and D configurations recur across these stages.
Identical inputs, seeds, objectives, and schedules define repeated executions of the same seeded fit.
Execution totals therefore do not increase the number of independent training seeds.
Each inference family uses only its declared seed and assignment sets.

New strengthening results require available prediction archives and completed analysis receipts.
Earlier reported strengthening outputs do not supply numbers for this revision.
The repository recovery ledger records that execution boundary.

\subsection{The frozen directional and retention design}
The retention design compares seven selected families for each encoder.
They are Source, Supported, allocation, distillation, WiSE-FT, image-source-only, and reverse-source-only.
The underlying grid has nine configurations, three rates, and three seeds.
This gives 81 base fits per encoder.
Nine matched randomized distillation controls raise each encoder's total to 90.
These controls copy D's primary selected coefficient, rate, and per-seed epoch.
They do not copy the later AD schedule.

Write $L_I^S,L_T^S$ for the source image and caption losses.
Write $L_I^U,L_T^U$ for the corresponding expanded losses.
The four directional cells are
\begin{align}
S&=\tfrac12(L_I^S+L_T^S),&
U&=\tfrac12(L_I^U+L_T^U),\\
I_S&=\tfrac12(L_I^S+L_T^U),&
T_S&=\tfrac12(L_I^U+L_T^S).
\end{align}
All four retain the same candidate captions.
$I_S$ restricts the image targets to source captions.
$T_S$ restricts eligible reverse anchors to source captions.
The reverse intervention changes which captions participate in the average.
It is not solely a change to individual image-target probabilities.

For outcome $Y$, the source-restriction main effects are
\begin{align}
\Delta_I&=\tfrac12[Y(S)+Y(I_S)-Y(U)-Y(T_S)],\\
\Delta_T&=\tfrac12[Y(S)+Y(T_S)-Y(U)-Y(I_S)].
\end{align}
The interaction is $Y(S)-Y(I_S)-Y(T_S)+Y(U)$.
These contrasts use epoch ten at each of the three rates.
Three contrasts, two retrieval endpoints, three rates, and two encoders give 36 effects.
Their 100,000 bootstrap resamples use seed 20261006.
The adjustment covers this complete directional family.

The separate selected-procedure family contains 80 contrasts.
Seven families against frozen contribute 28 effects per encoder across four endpoints.
D against A and WiSE-FT contributes eight.
D against matched randomized distillation contributes four.
The adjustment covers both encoders jointly.
Its 100,000 bootstrap resamples use seed 20261005.

Selection, coefficients, practical thresholds, and zero-loss sensitivity follow the unchanged frozen protocol.
The later AD extension adds its missing combined cell under a separate protocol.
Its factorial concerns allocation and distillation, not the image and reverse directional interventions.
All inference families remain distinct.

\subsection{Models and features}
ViT-B/32 uses the OpenCLIP LAION-2B checkpoint and its native score scale.
Its image and text features have 512 dimensions.
RN50 uses the official OpenAI checkpoint and its own preprocessing and score scale.
Its features have 1,024 dimensions.
Both encoders remain fixed during adaptation.

For a normalized feature $x$, a linear residual adapter returns
\begin{equation}
f_W(x)=\frac{x+Wx}{\|x+Wx\|_2}.
\end{equation}
The two modalities have independent matrices.
Both matrices start at zero.
The ViT adapter has 524,288 parameters; the RN50 adapter has 2,097,152 parameters.
No projection bias is trained.

The nonlinear ViT replication returns
\begin{equation}
f_{A,B}(x)=\frac{x+B\operatorname{ReLU}(Ax)}{\|x+B\operatorname{ReLU}(Ax)\|_2}.
\end{equation}
Each modality uses a 128-dimensional hidden layer.
The output matrix $B$ starts at zero.
The input matrix can learn after the output matrix changes.
The two branches contain 262,144 trainable parameters in total.
This replication changes adaptation capacity while retaining the same pretrained representation.

Every feature archive carries the encoder identity, checkpoint hash, preprocessing, image order, and text order.
Recomputed feature archives remain identified as fresh executions.
They are not described as exact-byte recoveries.

\subsection{Candidate and assignment invariants}
Training uses 1,200 images with their five source captions and all retained hypotheses.
Each hypothesis belongs to its annotated image.
Both replicated policies and random controls use the same complete batch candidates.
Unpromoted hypotheses remain competitors.

The score-stratified control ranks each image's hypotheses by its own encoder's frozen cosine scores.
Manifest index resolves exact score ties.
The ordered list splits into two nearly equal bins.
Each bin receives exactly its original supported count.
Sampling uses NumPy PCG64 seeds 101, 211, and 307.
Each assignment remains fixed across all rates, training seeds, epochs, and batches.
Naturally unchanged assignments remain in the sample.
The controls retain the original labels for evaluation and provenance.
Appendix~\ref{app:followupcontrols} gives the full assignment definition.

\subsection{Exact preservation objective}
Let $Z$ contain the scaled image--text scores for a batch.
Let $C$ contain only its source-caption columns.
Teacher and student row distributions are
\begin{align}
q_i^I&=\operatorname{softmax}(Z^0_{i,C}/T),\\
p_i^I&=\operatorname{softmax}(Z_{i,C}/T).
\end{align}
Column distributions $q_j^T,p_j^T$ normalize the same source submatrix over image rows.
The preservation term is
\begin{align}
D_S=\frac{1}{2B}\sum_{i=1}^{B}\KL(q_i^I\|p_i^I)
\nonumber\\
\quad+\frac{1}{2|C|}\sum_{j\in C}\KL(q_j^T\|p_j^T).
\label{eq:sourcekl}
\end{align}
Teacher probabilities do not receive gradients.
The temperature is $T=2$.
Equation~\ref{eq:adobjective} includes the multiplier $T^2$ outside this divergence.
Thus $\beta\in\{1,4,16\}$ gives effective multipliers 4, 16, and 64.

$L_S$ and $L_U$ each average image and eligible-caption losses with equal directional weight.
The mixture $\lambda L_S+(1-\lambda)L_U$ therefore mixes complete symmetric objectives.
It does not simply redistribute image-row targets while leaving the reverse loss unchanged.
For $K$ source captions and $M$ supports, the row source mass is
\begin{equation}
\lambda+(1-\lambda)\frac{K}{K+M}.
\end{equation}
The reverse source-anchor contribution also increases with $\lambda$.
Its precise mass depends on the current batch's eligible anchor counts.

\subsection{Complete fit grid}
Each encoder has 13 deterministic configurations.
Source and U contribute one each.
A contributes two allocation coefficients.
D contributes three distillation coefficients.
AD contributes all six coefficient pairs.
Three learning rates and three seeds produce 117 fits per encoder.
Nine matched randomized AD controls raise that count to 126.
The full extension therefore contains 252 fits.

Base fits use ten epochs, batches of 32 images, and the same seed-specific image order.
The rates are $10^{-4}$, $3\times10^{-4}$, and $10^{-3}$.
The seeds are 17, 29, and 43.
AdamW uses weight decay $.01$, coefficients $(.9,.999)$, and epsilon $10^{-8}$.
Gradient norms are clipped at one.
There is no learning-rate scheduler.
Deterministic CPU fitting uses two threads.
The full protocol records implementation and feature hashes before fitting.

WiSE-FT interpolates each Supported correction matrix with its zero initialization.
The candidate coefficients are $\alpha\in\{0,.1,.25,.5,.75,1\}$.
The same $\alpha$ applies to both modalities.
Interpolation uses model parameters before feature normalization.
It does not average normalized features or final scores.

\subsection{Selection rules}
The development retrieval pool has 900 images and 4,500 source captions.
It excludes the original 100 calibration images.
It includes the original 100-image validation subset used for relation accuracy.
This overlap is intentional and does not create a second independent development sample.

For each candidate and seed, an epoch is feasible when both development retrieval constraints hold.
The primary tolerance permits a one-percentage-point loss from each frozen reference.
Epoch zero is always available.
Among feasible epochs, selection maximizes image-balanced supported-versus-contradicted accuracy.
Relation ties receive half credit.
Exact ties retain the earlier epoch.

Each family then selects shared coefficients and a shared rate across the three seeds.
Selection maximizes the mean chosen relation score.
The protocol fixes the remaining tie order.
The zero-percentage-point tolerance repeats the procedure as a descriptive sensitivity analysis.
It does not introduce additional primary hypotheses.

Matched randomized AD controls copy the primary AD coefficients, rate, and per-seed epoch.
They use three assignment draws crossed with three training seeds.
They are fitted without an independent development search.
The fitted control run stops at the required selected epoch.
An epoch-zero match remains explicitly frozen.

\subsection{Independent selection and matched decomposition}
The primary comparisons evaluate independently selected Source, U, A, D, AD, and WiSE-FT procedures.
The matched random comparison evaluates caption identity under the primary AD schedule.
The decomposition instead takes all four cells at that AD-selected schedule.
It uses the selected allocation coefficient for A and AD.
It uses the selected distillation coefficient for D and AD.

For outcome $Y$, let $Y_{00},Y_{10},Y_{01},Y_{11}$ denote U, A, D, and AD respectively.
The conventional allocation effect is
\begin{equation}
\Delta_A=\tfrac12[(Y_{10}-Y_{00})+(Y_{11}-Y_{01})].
\end{equation}
The conventional distillation effect is
\begin{equation}
\Delta_D=\tfrac12[(Y_{01}-Y_{00})+(Y_{11}-Y_{10})].
\end{equation}
The interaction is
\begin{equation}
\Delta_{AD}=Y_{11}-Y_{01}-Y_{10}+Y_{00}.
\end{equation}
The interaction measures departure from additivity for this outcome and selected schedule.
It does not establish a universal mechanism across objectives or training distributions.

\subsection{Evaluation and inference}
Primary endpoints are e-ViL I$\to$T R@1, e-ViL T$\to$I R@1, relation accuracy, and SugarCrepe++ both-caption accuracy.
The preservation family contains 80 contrasts across two encoders.
The decomposition contains 24 contrasts: three factorial effects, four endpoints, and two encoders.
Each family uses 100,000 paired image-cluster bootstrap resamples and its declared Bonferroni adjustment.
The bootstrap retains each image's associated caption queries and benchmark records.
Randomized controls average assignment draws within each training seed before seed averaging.
The fixed seeds and draws are not themselves resampled.

The practical gate uses adjusted lower bounds against frozen initialization.
Both retrieval bounds must exceed $-1$ percentage point.
The SugarCrepe++ bound must exceed zero.
Each selected state must have a positive epoch and nonzero update norm.
An across-encoder success requires the same family to pass for both encoders.
Point estimates alone cannot establish the gate.
Retrieval intervals can cross zero if their lower bounds exceed the declared negative tolerance.

The replication families remain separate from preservation inference.
Each has 12 contrasts at epoch ten.
They compare Supported with Source and score-stratified promotion at every rate in both retrieval directions.
The original linear study retains its separate 18-contrast family.

Current strengthening predictions retain query identities and aggregate definitions.
Independent audits recompute reported effects from those predictions.
Source image overlap, checkpoint selection, and protocol hashes receive separate checks.
The complete tables below report every declared comparison.
